
\documentclass[12pt]{modelo/prova2}

\include{orientacoes}

\begin{document}
\maketitle

\question[0.5] Na arquitetura da Internet, os dispositivos dos usuários (como computadores e celulares) são chamados de:
\begin{choices}
\item Núcleo da rede
\item Sistemas finais (hosts)
\item Meios de transmissão
\end{choices}
\answer{B}

\question[0.5] Qual das opções abaixo é uma vantagem do cabo de fibra óptica em relação ao cabo UTP?
\begin{choices}
\item Menor custo de instalação
\item Maior imunidade a interferência eletromagnética
\item Mais fácil de conectar em switches comuns
\end{choices}
\answer{B}

\question[0.5] Em uma topologia estrela, os dispositivos são conectados a:
\begin{choices}
\item Um dispositivo central (como um switch)
\item Um único cabo compartilhado por todos
\item Vários cabos formando um anel fechado
\end{choices}
\answer{A}

\question[0.5] O modelo OSI divide a comunicação em redes em:
\begin{choices}
\item 4 camadas
\item 5 camadas
\item 7 camadas
\end{choices}
\answer{C}


\question[0.5] Na arquitetura TCP/IP, o protocolo IP é responsável por:
\begin{choices}
\item Garantir que os dados cheguem corretamente ao destino
\item Encaminhar os pacotes entre redes diferentes (roteamento)
\item Criptografar os dados antes da transmissão
\end{choices}
\answer{B}


\question[0.5] O padrão Ethernet (IEEE 802.3) é mais comumente utilizado em:
\begin{choices}
\item Redes sem fio (Wi-Fi)
\item Redes com fio (cabeadas)
\item Redes de telefonia celular
\end{choices}
\answer{B}


\question[0.5] Um switch de camada 2 utiliza qual informação para encaminhar os quadros?
\begin{choices}
\item Endereço IP de destino
\item Endereço MAC de destino
\item Nome do dispositivo
\end{choices}
\answer{B}


\question[0.5] O cabo de par trançado (UTP) é composto por:
\begin{choices}
\item 2 pares de fios
\item 4 pares de fios
\item 8 pares de fios
\end{choices}
\answer{B}

\question[0.5] O padrão IEEE 802.11 é utilizado em:
\begin{choices}
\item Redes com fio (Ethernet)
\item Redes sem fio (Wi-Fi)
\item Redes de fibra óptica
\end{choices}
\answer{B}


\question[0.5] Em uma rede Ethernet, um hub (concentrador):
\begin{choices}
\item Segmenta os domínios de colisão
\item Compartilha o mesmo domínio de colisão entre todas as portas
\item Elimina completamente as colisões
\end{choices}
\answer{B}

\noindent\textbf{INSTRUÇÕES PARA AS QUESTÕES DISCURSIVAS:}
\begin{itemize}
\item Responda com frases curtas
\item Você pode usar desenhos ou esquemas se preferir
\item Não se preocupe com o tamanho da resposta, o importante é mostrar que você entendeu
\end{itemize}

\question[1] Um hospital precisa escolher entre cabo UTP e fibra óptica para sua rede.

A) Qual meio de transmissão é mais adequado para ambientes com muita interferência eletromagnética, como um hospital? Por quê?

\vspace{5\baselineskip}

B) Qual meio de transmissão costuma ter menor custo para curtas distâncias?

\vspace{5\baselineskip}


\question[2] Um professor montou uma rede com 10 computadores usando um hub. Os alunos reclamaram que a rede está muito lenta.

A) Explique por que um hub pode deixar a rede lenta quando muitos computadores estão conectados.

\vspace{5\baselineskip}

B) Que equipamento poderia substituir o hub para melhorar o desempenho da rede?

\vspace{5\baselineskip}


\question[2] Sobre rede cabeada e sem fio, responda:

A) Em uma situação em que a estabilidade da conexão é essencial, como em um servidor de arquivos de uma empresa, qual tipo de rede é mais recomendado? Por quê?

\vspace{5\baselineskip}

B) Cite uma vantagem da rede sem fio (Wi-Fi) em relação à rede cabeada.

\vspace{5\baselineskip}

\end{document}
